\documentclass[11pt]{article}
\usepackage[a4paper,margin=26mm]{geometry}
\usepackage{amsmath,amssymb,amsthm,hyperref}
\setlength{\emergencystretch}{3em}
\newtheorem{theorem}{Theorem}
\newtheorem{lemma}[theorem]{Lemma}
\newtheorem{proposition}[theorem]{Proposition}
\newtheorem{corollary}[theorem]{Corollary}
\newcommand{\PSH}{\operatorname{PSH}}
\newcommand{\MI}{\mathcal I}
\newcommand{\OO}{\mathcal O}
\newcommand{\NA}{\mathrm{NA}}
\newcommand{\MA}{\operatorname{MA}}
\newcommand{\ddc}{dd^c}
\newcommand{\ds}{d_{\mathcal S}}
\newcommand{\Eone}{\mathcal E^1}
\title{A transcendental obstruction to model approximation\protect\\ preserving the finer singularity profile}
\author{Unreviewed original P19.2 counterexample candidate, continuation 211}
\date{6 October 2026}
\begin{document}
\maketitle
\begin{abstract}
We construct a bounded finite-energy weak geodesic ray from the smooth
zero potential on a nonprojective compact K\"ahler threefold with a
genuinely transcendental K\"ahler class. Its usual non-Archimedean datum
is zero, but every interior Legendre level in a nonempty interval has a
positive-mass P-model singularity that is not I-model. No sequence of
canonical cohomological model rays, or even of arbitrary usual-datum
maximal finite-energy rays from zero, converges to this finer profile
in the singularity distance. The obstruction already concerns positive
mass and is independent of any added entropy requirement. We prove the
needed I-model closedness using lower rooftop envelopes on an arbitrary
compact K\"ahler manifold, and justify the local Gauss formula without
assuming a polarization. This is a proposed analytic counterexample to
the arbitrary-ray, finer-P-profile approximation request. It remains
unreviewed; no original problem is recorded as solved on its basis.
\end{abstract}

\section{The approximation signature and conventions}
Let $(X,\omega)$ be a connected smooth compact K\"ahler manifold,
$n=\dim X$, $V=\int_X\omega^n$, and $\alpha=[\omega]$. We fix
\[
 \ddc=\frac{i}{2\pi}\partial\bar\partial,
 \qquad t=-\log|\tau|^2,
 \qquad \ddc\log|f|^2=[\operatorname{div}(f)].                 \tag{1}
\]
All products of unbounded potentials below are nonpluripolar products.
Set
\[
 M_k(u)=\frac1V\int_X\langle(\omega+\ddc u)^k\rangle
                      \wedge\omega^{n-k},\quad 0\le k\le n.
                                                               \tag{2}
\]
Thus $M_0=1$ and $M_n$ is a raw mass divided by the fixed reference
volume, not a renormalized probability. Put
\[
 E(u)=\frac1{(n+1)V}\sum_{k=0}^n
     \int_X u\,\omega_u^k\wedge\omega^{n-k},\qquad
 H(u)=\operatorname{Ent}(\omega_u^n/V\mid\omega^n/V)
                                                               \tag{3}
\]
on the full-mass finite-energy class, with entropy $+\infty$ for a
measure not absolutely continuous with respect to $\omega^n/V$.

Write $P[u]=\operatorname{usc}\lim_{C\to\infty}P_\omega(u+C,0)$.
The I-envelope is
\[
 P_I[u]=\sup\{v\in\PSH(X,\omega):v\le0,
                 \MI(av)\subset\MI(au)\text{ for every }a>0\}.
                                                               \tag{4}
\]
The envelope in (4) is itself admissible. A normalized P-model is a
fixed point of $P$, and an I-model is a fixed point of $P_I$.
Positive-mass P-models carry the singularity distance of DDNL
\cite{DDNL}. We use its probability normalization $\ds^V=\ds/V$.
In particular, for ordered types $[u]\le[v]$,
\[
 \ds^V([u],[v])=\frac1{n+1}\sum_{k=0}^n(M_k(v)-M_k(u)).        \tag{5}
\]
Here P-equivalence refers to equality of the normalized P-envelopes;
it is finer than equality of all multiplier ideals.

For an actual finite-energy weak ray $U$ from zero, let
\[
 \chi_\lambda=\inf_{t>0}(U_t-t\lambda),\qquad
 U_t=\sup_\lambda(\chi_\lambda+t\lambda).                    \tag{6}
\]
These are the closed Legendre levels; a level identically $-\infty$
is not a potential to which (2) or a positive-mass singularity distance
can be assigned. The finer profile convergence examined here requires
\[
 \ds^V([P[\chi_{j,\lambda}]],[\chi_\lambda])\longrightarrow0
                                                               \tag{7}
\]
at every genuine positive-mass level of the target. In particular (7)
preserves the finer P datum and the raw mixed masses in (2).
The prescribed start is zero, so these normalized absolute envelopes
are also the relative P ceilings against the actual start.
The counterexample fails (7) already at any one of the indicated
interior levels, so additional requirements at a closed endpoint do
not repair it.

A cohomological model means a dominating equivariant normal compact
K\"ahler test configuration, with a real vertical divisor $D$, whose
class $\alpha_{\mathcal X}+D$ is relatively K\"ahler, together with its
canonical Perron weak ray from zero. We also allow canonical maximal
rays attached to finite-energy nef or obstacle model data. This includes
positive analytic flag blowup configurations; no relative semiampleness
or line bundle is imposed. More broadly, our obstruction applies to
\emph{every} usual-datum maximal finite-energy ray from zero.
Maximality means domination of any weak ray with smaller initial
potential and smaller usual datum. The usual datum is the negative
generic coefficient of the actual total potential on a vertical Gauss
divisor, divided by its order of $\tau$, in convention (1).
It is extended to $X^\beth$ by divisorial determination, retaining
supported semivaluations. These notions and the actual maximal rays
are the complete previously accepted caller theorem R \cite{R};
the relevant geometric model realization is MP Lemma 5.1.8 \cite{MP}.

\begin{theorem}\label{main}
There are a nonprojective compact K\"ahler threefold $(X,\omega)$,
with $[\omega]$ not proportional to a line-bundle class, and an actual
weak geodesic ray $U:[0,\infty)\to\Eone(X,\omega)$ such that
\begin{enumerate}
\item $U_0=0$, $-t\le U_t\le0$, and its usual datum is identically zero.
\item For a fixed $c_0\in(0,1)$, its Legendre levels are P-models and
\[
 M_k(\chi_\lambda)=1-\frac{k}{n}c_0(\lambda+1),
       \qquad -1<\lambda\le0,
       \quad 0\le k\le n.                                 \tag{8}
\]
Every level in $(-1,0]$ has all multiplier ideals trivial and has
positive mass strictly less than one. Hence it is not I-model.
\item For any fixed $\lambda\in(-1,0)$, no sequence of usual-datum
maximal finite-energy rays from zero satisfies (7). Consequently no
cohomological model approximation of this ray preserves this finer
P profile, regardless of the energy and entropy conditions added.
\item Its radial energy is $-c_0/2$, whereas the maximal ray of its
usual datum from zero is stationary and has radial energy zero.
With the entropy normalization (3), $H(U_0)=0$ and
$\lim_{t\to\infty}H(U_t)/t=+\infty$ by caller theorem R.
\end{enumerate}
\end{theorem}

\section{Closedness of I-model singularities without a polarization}
We specify the external envelope facts used. DX Proposition 2.18 and
Lemma 2.21 are valid in its Section 2.4 on an arbitrary compact
K\"ahler manifold, with a smooth pseudoeffective form: $P_Iu$ is
P-model, $P_Iu\sim_Iu$, and decreasing I-model limits remain I-model.
An increasing almost-everywhere limit of I-model P-potentials is
I-model when the limit is P-model of positive mass \cite{DX}.
DDNL Theorem 5.6 applies to a big class on a compact K\"ahler
manifold: if positive-mass normalized P-models converge in $\ds$,
a subsequence has lower rooftop tails increasing almost everywhere
to the limit. The tails are decreasing limits of finite rooftops,
are P-model, and have positive mass after a finite initial deletion
\cite{DDNL}. We need the lower tails, which require no algebraic
approximation.

\begin{lemma}\label{roof}
A nontrivial finite lower rooftop of I-model potentials is I-model.
\end{lemma}
\begin{proof}
Let $w=P_\omega(u_1,\ldots,u_r)\in\PSH(X,\omega)$, where
$P_Iu_i=u_i$. Since $w\le u_i$, any competitor $v$ in (4) for $w$
satisfies $\MI(av)\subset\MI(aw)\subset\MI(au_i)$ for all $a>0$.
Thus $v\le P_Iu_i=u_i$ for every $i$, and $v\le w$ by the
rooftop definition. The potential $w$ itself is a competitor.
Taking the supremum proves $P_Iw=w$.
\end{proof}

\begin{proposition}\label{closed}
If $u_j=P_Iu_j$, $u=P[u]$, $M_n(u)>0$, and
$\ds([u_j],[u])\to0$, then $u=P_Iu$.
\end{proposition}
\begin{proof}
The mass continuity bound following DDNL Lemma 3.4, Proposition 3.5
and Lemma 3.7 gives
$|M_n(u_j)-M_n(u)|\le C\ds^V([u_j],[u])$; it follows directly by
splitting at $\max(u_j,u)$ and using the nonnegative summands of
(5). Thus the sequence and its limit eventually have a common
positive lower mass bound. Apply DDNL Theorem 5.6 and relabel its
subsequence $u_j$. The potentials
\[
 w_j=P_\omega(u_j,u_{j+1},\ldots)
       =\lim_{r\to\infty}P_\omega(u_j,\ldots,u_{j+r})
                                                               \tag{9}
\]
are P-models and increase almost everywhere to $u$. In the theorem's
construction all the finite rooftops in (9) are nontrivial, and the
tails are nontrivial of positive mass, after deleting finitely many
indices. Lemma \ref{roof} makes every finite rooftop I-model.
DX Lemma 2.21(i) makes $w_j$ I-model on passage to the decreasing
limit in $r$. Its Lemma 2.21(iii), now applied to $w_j\nearrow u$
and using $M_n(u)>0$, gives $P_Iu=u$. All applications concern the
same fixed K\"ahler class; neither global sections nor a polarization
are present.
\end{proof}

\section{Why usual maximal rays have I-model interior levels}
The next argument supplies the nonprojective step explicitly, rather
than importing the projective ambient hypothesis in DXZ Section 3.1.

\begin{lemma}[Local Gauss formula]\label{gauss}
Let $\chi_\lambda$ be a test curve from zero with cutoff at most zero,
and let $W_t=\sup_\lambda(\chi_\lambda+t\lambda)$ be its sublinear
subgeodesic. For a rationally scaled divisorial valuation $v$ over a compact complex
manifold, let $\nu_v(\chi)$ be its generic divisorial coefficient in
convention (1). Then
\[
 W^\NA(v)=\sup_{\lambda<\lambda_+}
                          (\lambda-\nu_v(\chi_\lambda)).    \tag{10}
\]
The formula holds for the analytic Gauss divisors used in $X^\beth$;
it does not require that $X$ be projective.
\end{lemma}
\begin{proof}
On $X\times\Delta$ the total potential is the regularized supremum
of $\chi_\lambda(x)-\lambda\log|\tau|^2$. In a coordinate chart,
adding a local smooth potential of $\omega$ makes these functions
psh. They are uniformly bounded above, since $\chi_\lambda\le0$,
$\lambda\le0$ and $|\tau|<1$. The same is true on a modification.
For a locally uniformly upper bounded family of psh functions,
the Lelong number of its regularized supremum equals the infimum of
the individual Lelong numbers. Indeed Choquet's lemma reduces to a
countable family, finite maxima have the minimum Lelong number, and
their increasing $L^1$ limit has that same limiting number by upper
semicontinuity and monotonicity. This is the local proof of DX
Lemma 3.14, whose ambient space is any complex manifold.

Let $F$ be a normalized Gauss divisor with $b=\operatorname{ord}_F\tau$.
The coefficient of the summand is
$b(\nu_v(\chi_\lambda)-\lambda)$. This is a local coefficient
calculation: after making the divisor underlying $v$ smooth, the
Gauss valuation is realized by the weighted blowup in its normal
coordinate and $\tau$. The pullback coefficient of a potential
independent of $\tau$ is its normal-coordinate coefficient times
the normal weight; the coefficient of $\log|\tau|^2$ is $b$.
Equivalently one uses the analytic ideal Gauss rule
$\sigma(v)(\sum a_i\tau^i)=\min_i(v(a_i)+i)$.
Analytic Gauss divisors and their restrictions are precisely those
of MP Lemma 2.4.2 and Theorem 2.5.5. Apply the supremum rule at
generic points of $F$, divide by $b$, and negate the coefficient.
The cutoff can be approached from the left by concavity; the
Legendre supremum is unchanged if only $\lambda<\lambda_+$ is used.
This proves (10) with the normalization of the actual usual datum.
\end{proof}

\begin{proposition}\label{maxI}
For any usual-datum maximal finite-energy ray $U$ from zero, every
genuine interior Legendre level $\chi_\lambda$, $\lambda<\lambda_+$,
is I-model. In particular this holds for the canonical model rays
in Section 1.
\end{proposition}
\begin{proof}
Subtract an affine time function if necessary to put the cutoff at
most zero. Legendre duality on a compact K\"ahler manifold makes
the levels normalized P-models. Apply the I-maximization of test
curves in DXZ Section 2.2: replace interior levels by
$\eta_\lambda=P_I[\chi_\lambda]$, take the left limit at the cutoff,
and put $-\infty$ above it. This is an I-maximal test curve from
$P_I[0]=0$. Its inverse Legendre transform $W$ is a weak geodesic
ray from zero and satisfies $W_t\ge U_t$. It has finite energy:
each level has no smaller mass, and the finite-energy test-curve
integral $\int_{-\infty}^{\lambda_+}(M_n-1)\,d\lambda$ is
increased and is still at most zero.

By DX Proposition 2.18(ii) the interior old and new levels have the
same multiplier ideals at every positive coefficient. DX Theorem
2.13, in its general compact-K\"ahler Section 2.4, identifies their
coefficients along all projective analytic modifications. Thus
Lemma \ref{gauss} gives $W^\NA=U^\NA$ on divisorial points. The
usual psh extensions agree on $X^\beth$ by caller R's divisorial
determination. Maximality of $U$, with the same initial zero, now
gives $U\ge W$. Hence $U=W$. Injectivity of the Legendre
correspondence gives $\chi_\lambda=P_I\chi_\lambda$ at every
interior level. Restoring the affine time shift merely translates
the level parameter and does not change this assertion.
\end{proof}

\section{An explicit diffuse polar singularity}
We construct the measure to ensure polar support; a general nonatomic
measure on a Cantor set would not suffice. Inside a small real interval
in the affine chart of $C=\mathbf P^1$, form nested binary intervals.
At stage $m$ each of the $2^m$ intervals has radius
$r_m=\exp(-4^m)$; its two children have radius $r_{m+1}$ and centers
near the two quarter points of the parent. They are disjoint and
contained in it. Let $K$ be the intersection of their unions and let
$\nu$ be the Bernoulli measure giving each stage-$m$ interval mass
$2^{-m}$. It is a nonatomic probability supported on $K$.
Use the Fubini--Study form $\omega_C$ of mass one and its Green kernel
\[
 G(z,a)=\log|z-a|^2-\log(1+|z|^2)-\log(1+|a|^2)\le0,
 \qquad h(z)=\int_K G(z,a)\,d\nu(a).                         \tag{11}
\]
The kernel is understood in homogeneous coordinates at infinity.
Its negative part has a uniformly bounded integral against smooth
area, so $h\in L^1(C)$ by Fubini. Also
$\ddc h=\nu-\omega_C$, so $h$ is $\omega_C$-psh.
If $z\in K$, the cylinder containing $z$ gives
\[
 h(z)\le 2^{-m}\log(4r_m^2)\longrightarrow-\infty.          \tag{12}
\]
Outside $K$, the kernel is uniformly finite on the support of $\nu$
locally, so $h$ is finite. Thus its pole set is exactly $K$, which
in particular is polar and has smooth area zero.

\begin{lemma}\label{integrability}
For every $a>0$, $e^{-a h}$ is locally integrable. Consequently,
for $0<c<1$, $u_c=ch$ has all multiplier ideals trivial and
its nonpluripolar mass is $1-c$.
\end{lemma}
\begin{proof}
Atomlessness on the compact set implies
$\sup_x\nu(B(x,r))\to0$ as $r\to0$: otherwise choose shrinking
balls of a fixed positive mass, pass their centers to a limit, and
obtain an atom. Fix $a>0$. Around any point choose a ball whose
doubled ball has mass $m<1/a$. For points in the smaller ball,
the far part of (11) is bounded below. If $m>0$, Jensen's inequality
for the probability $m^{-1}\nu$ on the doubled ball gives
\[
 e^{-a h(z)}\le C\int |z-b|^{-2am}\,\frac{d\nu(b)}m.
                                                               \tag{13}
\]
The smooth terms in the Green kernel are included in $C$. Because
$am<1$, integrate (13) over the smaller ball and apply Fubini;
the integrals of $|z-b|^{-2am}$ are uniformly finite. If $m=0$
the potential is locally bounded there. The argument applies to
any exponent and hence also to $ch$.

Finally $\omega_C+\ddc(ch)=(1-c)\omega_C+c\nu$.
In dimension one its nonpluripolar part is the restriction to the
complement of the pole set. The part $c\nu$ is removed, while
$\omega_C(K)=0$, giving mass $1-c$.
\end{proof}

Let $q_c=P_C[ch]$. The P-envelope preserves top mass and is model
at positive mass, so
\[
 q_c\ge ch,\quad q_c\le0,\quad
 \int_C\langle\omega_C+\ddc q_c\rangle=1-c,
 \quad \MI(aq_c)=\OO_C\ (a>0).                             \tag{14}
\]
The mass-preserving envelope assertion and its positive-mass
idempotence are DDNL Section 2 and Proposition 2.6, in their big
class scope; here the class is K\"ahler. Integrability in (14)
also follows directly from $q_c\ge ch$ and Lemma \ref{integrability}.
Thus $P_Iq_c=0$, whereas $q_c\ne0$ by its mass. The diffuse-polar
mechanism on a projective curve is the known distinction in DX
Example 2.19; our extension below is on a nonprojective manifold,
with an explicit actual ray and a general analytic closedness proof.

\section{The nonprojective manifold and its product envelopes}
Choose a two-dimensional complex torus $Y$ with
$H^{1,1}(Y)\cap H^2(Y,\mathbf Z)=\{0\}$. For completeness such
a torus exists by the following elementary construction. On
$\mathbf R^4/\mathbf Z^4$ consider the connected space of complex
structures with a fixed orientation. For each nonzero integral
alternating form $A$, the condition that $A$ be J-invariant is a
proper closed real analytic condition. To see properness choose an
oriented basis with $A(e_1,e_2)\ne0$, and put
\[
 Je_1=s e_3,\quad Je_3=-s^{-1}e_1,\quad
 Je_2=t e_4,\quad Je_4=-t^{-1}e_2.
\]
Take $st<0$ for the chosen complex orientation, and vary $st$;
$A(Je_1,Je_2)=stA(e_3,e_4)$ cannot always equal $A(e_1,e_2)$.
Baire's theorem avoids all these countably many conditions.
The averaging description of torus cohomology identifies the
integral $(1,1)$ classes with these invariant forms. Any constant
positive Hermitian form is K\"ahler. An ample bundle would have a
nonzero integral positive $(1,1)$ class, so $Y$ is not projective.

Put
\[
 X=C\times Y,\qquad \omega=p_C^*\omega_C+p_Y^*\omega_Y,
 \qquad n=3,\qquad V=n\int_Y\omega_Y^{n-1}.                 \tag{15}
\]
The same argument works with any nonprojective K\"ahler factor
having no nonzero integral $(1,1)$ class. A projective $X$ would
make its closed fiber $Y$ projective. Moreover $[\omega]$ cannot
be proportional to a line-bundle class: restrict that class to $Y$.

\begin{lemma}\label{product}
For $f\in\PSH(C,\omega_C)$, $f\le0$,
\[
 P_X[p_C^*f]=p_C^*P_C[f].                                  \tag{16}
\]
For $Q_c=p_C^*q_c$ one has
\[
 P_X[Q_c]=Q_c,\quad P_{I,X}[Q_c]=0,\quad
 M_k(Q_c)=1-\frac{k c}{n}\quad(0\le k\le n).                \tag{17}
\]
\end{lemma}
\begin{proof}
First work with the finite rooftop $P_X(p_C^*f+A,0)$.
Pulling back a curve competitor gives a product competitor.
Conversely, restricting a product competitor to each fiber
$C\times\{y\}$ gives an $\omega_C$-psh function or the identically
$-\infty$ function, bounded by $\min(f+A,0)$. It is therefore at
most $P_C(f+A,0)$. This upper bound is usc on the product, so it
also bounds the regularized supremum of the competitors. This
proves equality for every $A$, and passage to the increasing
regularized limit proves (16).

The multiplier ideals of $Q_c$ are trivial by the local product
Fubini theorem and (14). Thus zero is admissible in (4) and its
I-envelope is zero. To calculate the products, truncate $q_c$ by
$\max(q_c,-N)$ on the curve. The resulting current on the product
is the sum of its curve current and $p_Y^*\omega_Y$; two curve
factors vanish. After restriction to $\{q_c>-N\}$ and passage
to the nonpluripolar limit, the integrand in (2) is
\[
 \big((n-k)\omega_C+k\langle\omega_C+\ddc q_c\rangle\big)
                       \wedge\omega_Y^{n-1}.               \tag{18}
\]
The pole set of $q_c$ has smooth area zero, so no smooth curve
mass is lost in the first term. Its integral is
$(n-k)+k(1-c)=n-kc$ times $\int_Y\omega_Y^{n-1}$.
Divide by (15) to obtain (17). In particular these are positive
raw masses, and none is replaced by a probability measure at a
singular level.
\end{proof}

\section{The actual weak ray and failure of every model sequence}
Fix $0<c_0<1$ and start with the raw test curve on $C$
\[
 a_\lambda=\begin{cases}
 0,&\lambda\le-1,\\
 c_0(\lambda+1)h,&-1<\lambda\le0,\\
 -\infty,&\lambda>0.
 \end{cases}                                               \tag{19}
\]
It is decreasing, concave and upper semicontinuous in the level
parameter, including the extended values at points of $K$.
Its limit at $-\infty$ is zero. P-maximize (19) as in DXZ
Section 2.2 and formula (13). For $-1<\lambda<0$ its levels are
$q_{c_0(\lambda+1)}$, below $-1$ they are zero, and the value at
zero is their decreasing left limit. That endpoint equals $q_{c_0}$:
DDNL Proposition 4.8 and Corollary 4.7 make the left limit P-model
of mass $1-c_0$; it dominates $q_{c_0}$ with the same positive
mass. DDNL Theorem 2.3 applied to this ordered pair makes their
P-envelopes equal. Pull the maximized curve back to $X$.
Lemma \ref{product} shows that it is still a maximal test curve;
its concavity and closed endpoints are preserved by pullback.
Denote its levels by $\chi_\lambda$.

The relative Ross--Witt Nystr\"om theorem, DXZ Theorem 2.6(ii)--(iii),
applies to this compact K\"ahler manifold and positive-volume
reference zero. Hence the inverse transform (6) is an actual
finite-energy weak geodesic ray. More concretely the available
level $\lambda=-1$ and the upper bound on all levels give
$-t\le U_t\le0$. Thus all times are bounded finite-energy
potentials, and $U_t\to0$ strongly as $t\downarrow0$.
Its spectrum is contained in $[-1,0]$. The energy formula in that
same theorem and (17) give
\[
 \lim_{t\to\infty}\frac{E(U_t)}t
 =\int_{-1}^0\big(M_n(\chi_\lambda)-1\big)\,d\lambda
 =-\int_{-1}^0c_0(\lambda+1)\,d\lambda=-\frac{c_0}2.        \tag{20}
\]
The explicit mixed masses in (17) prove (8), including the genuine
closed endpoint at zero. Formula (5) also gives the useful fixed
scale comparison
\[
 \ds^V([\chi_\lambda],[0])=\tfrac12c_0(\lambda+1)>0
                    \quad(-1<\lambda\le0).                \tag{21}
\]

All genuine levels have trivial multiplier ideals. They therefore
have zero generic coefficients on all projective analytic
modifications by DX Theorem 2.13, applied to comparison with zero.
Lemma \ref{gauss} gives $U^\NA(v)=\sup_{\lambda<0}\lambda=0$
on divisorial points, and caller R extends it as the zero function
on the full carrier. Its maximal ray from zero is the stationary
ray. Equation (20) proves that $U$ is not that ray. Caller R's
arbitrary-ray entropy classification, with the initial entropy
$H(0)=0$, gives its infinite slope asserted in Theorem \ref{main}.

Now fix $\lambda\in(-1,0)$ and suppose a sequence $U^j$ of
usual-datum maximal rays from zero satisfied (7). A sequence that
is identically $-\infty$ at this level infinitely often already
fails the genuine-level requirement. For the remaining sequence,
Proposition \ref{maxI} makes each interior level I-model. If a
level happens to be the cutoff of an approximating ray, the same
holds: its left limits are I-model and the decreasing-limit
part of DX Lemma 2.21 applies whenever the cutoff level is
genuine. Thus $P[\chi_{j,\lambda}]$ is I-model.
Since the target level has mass $1-c_0(\lambda+1)>0$,
Proposition \ref{closed} would make it I-model. But its
I-envelope is zero while its mass is strictly less than one,
a contradiction. This proves all the claims of Theorem \ref{main}.

\section{Comparison with the original residual problem}
Original P19.2 asks for model approximations preserving required
energy, entropy and singularities for arbitrary finite-energy weak
rays on compact K\"ahler manifolds; it explicitly keeps finer P
data separate from I data. The signature (7) is the usual singularity
distance on its positive-mass P-model levels, and automatically
records every mixed mass in (2). Our obstruction is at an interior
positive-mass level of a bounded actual ray from a smooth finite-entropy
start, on a nonprojective manifold in a genuinely transcendental
class. It therefore does not rely on absolute/relative changes at a
singular start, raw zero mass, an identically $-\infty$ placeholder,
or an erroneous choice of one approximation sequence.

Propositions \ref{maxI} and \ref{closed} exclude every canonical
model ray in the stated cohomological category, and in fact a larger
category of all usual-datum maximal rays. The result rules out the
unrestricted arbitrary-ray approximation request when its finer
P singularities are required in this signature. Energy or entropy
convergence cannot supply the missing P data. It leaves possible
approximation statements for I-maximal rays, for coarser ideal data,
or using additional genuinely nonmodel complex singularities as
input. The known curve distinction alone would not settle a residual
transcendental problem; here the general analytic closedness argument,
the nonprojective product and the actual ray supply that ambient
extension. P19.3 and P19.4 concern different goals and are unaffected.

This manuscript is a proposed full-original counterexample for this
frozen finer-profile signature, conditional only on the complete
accepted caller theorem R and the precisely cited ordinary envelope
and Legendre results. Independent review must check both its proof
and whether this signature answers the exact original completion
criterion. There is currently no independent verdict, and original
P19.2 remains open in the research catalogue.

\appendix
\section{Source application and reading boundaries}
The complete accepted caller R is reused as an explicit mathematical
grant; it supplies the actual usual-datum carrier, divisorial
determination, maximality and the entropy classification. Its new
applications here use the fixed smooth compact K\"ahler class,
start zero, and convention (1). No theorem from its future
P19.2/P19.3/P19.4 programme is imported.

Primary versions were downloaded on 6 October 2026 to a private
reading directory and bound by hashes in a separate application
record. The pertinent complete units read are DX's general
Section 2.4 definitions, Proposition 2.18 and Lemma 2.21, and
its local Lemma 3.14; DDNL's Lemma 4.3, Propositions 4.6 and 4.8,
Corollary 4.7, Lemmas 5.1--5.2 and full Theorems 5.4 and 5.6;
and DXZ's relative
Section 2.2, Legendre theorem and maximization constructions.
The cited mass preservation and full-mass characterization in
DDNL Section 2 were read with their stated hypotheses. MP's
model lemma and usual maximality statements were read with their
proofs; its analytic Gauss theorem statement and local reduction
were read, but its whole recursive local-algebraic proof was not
reconstructed here. That carrier is already part of caller R.
No claim is made to have read an entire primary paper or to have
independently reconstructed every pluripotential foundation.
The published projective-only hypothesis of DXZ Section 3.1 is
retained: Lemma \ref{gauss} gives the local extension needed here.

\begin{thebibliography}{9}
\bibitem{R} Complete accepted caller theorem R,
\emph{Entropy slopes of finite-energy K\"ahler rays: calibration,
recovery, and arbitrary starting points}, original P19.1 paper,
6 October 2026; supplied in full to the independent review.
\bibitem{DX} T. Darvas and M. Xia, \emph{The closures of test
configurations and algebraic singularity types}, Advances in
Mathematics 397 (2022), 108198,
\href{https://research.chalmers.se/publication/528300/file/528300_Fulltext.pdf}{published primary PDF}.
\bibitem{DDNL} T. Darvas, E. Di Nezza and C. H. Lu,
\emph{The metric geometry of singularity types},
\href{https://arxiv.org/pdf/1909.00839v1}{arXiv:1909.00839v1},
2 September 2019, Sections 2--5.
\bibitem{DXZ} T. Darvas, M. Xia and K. Zhang,
\emph{A transcendental approach to non-Archimedean metrics of
pseudoeffective classes}, Commentarii Mathematici Helvetici 100
(2025), 269--322,
\href{https://ems.press/content/serial-article-files/49836}{published primary PDF}, Section 2.2.
\bibitem{MP} P. Mesquita-Piccione,
\emph{A non-Archimedean theory of complex spaces and the cscK problem},
\href{https://arxiv.org/pdf/2409.06221v2}{arXiv:2409.06221v2},
Sections 2.4--2.5 and 5.1; the same fixed-manifold carrier as caller R.
\end{thebibliography}
\end{document}
